% =============================================================================
% Chapter 12 -- Causality Study: Motivation and Design
% =============================================================================

\chapter{Causality Study: Motivation and Design}
\label{ch:causality_intro}

Part~IV showed that circular oversampling improves minority-class classification, but all such work rests on an assumption that has never been directly tested: that class imbalance \emph{per se} causes the degradation.
Most empirical studies compare oversamplers against each other on datasets that are already imbalanced at a fixed, unknown ratio and do not control for confounds: reduced minority sample size, distributional shift from naive removal, or classifier-family sensitivity.
Part~V tests this causal assumption directly.
A uniform random removal protocol drives the degradation phase; GVM-CO (Chapter~\ref{ch:proposed_methods}) is one of the recovery oversamplers evaluated.


% ============================================================================
\section{The Untested Assumption}
\label{sec:causal:assumption}

Machine learning practitioners routinely encounter a well-known heuristic: when a binary classifier performs poorly on the minority class, class imbalance is blamed.
This assumption is so pervasive that entire research sub-fields have emerged around oversampling, undersampling, and cost-sensitive learning~\cite{he2009learning, fernandez2018learning}, all premised on the idea that an unequal class distribution \emph{causes} degraded classification performance.

Yet the causal nature of this claim has rarely been tested with adequate rigour.
Most empirical studies compare different oversampling techniques against one another, or against a no-oversampling baseline, but do so on datasets that are \emph{already} imbalanced at an unknown and fixed ratio~\cite{krawczyk2016learning, haixiang2017learning}.
They do not isolate the effect of imbalance itself from the confounding effects listed above.

As a result, while the literature documents \emph{that} classifiers struggle under imbalance, the question of \emph{why} --- and specifically whether the ratio itself or the accompanying information loss is the root cause --- remains underexplored.
Japkowicz and Stephen~\cite{japkowicz2002class} raised this question early, noting that class imbalance interacts with class overlap in non-obvious ways, but a controlled causal study was not performed.

This part of the thesis addresses the gap directly.
We ask: \emph{Does progressively increasing class imbalance monotonically degrade classifier performance, and can oversampling reverse that degradation to recover the balanced baseline?}


% ============================================================================
\section{Study Design}
\label{sec:causal:design}

We design a controlled \emph{degradation-and-recovery} experiment on ten binary classification datasets that begin in a fully balanced state (imbalance ratio $\IR \approx 1$).
The experiment proceeds in two phases.

\textbf{Phase~1 --- Degradation sweep.}
Starting from the balanced dataset, we remove minority-class instances one step at a time (1\% of the original minority per step) using a uniform random removal protocol (Section~\ref{sec:causal:random_removal}).
We continue until 80\% of the minority class has been removed ($\IR \approx 4{:}1$).
At each step we measure six classification metrics across five classifiers using leakage-safe five-fold cross-validation (Section~\ref{sec:setup:leakage}).

\textbf{Phase~2 --- Recovery sweep.}
Starting from the most-degraded state, we apply each of seven oversampling methods incrementally, adding back 1\% of minority samples per step until full balance is restored.
We measure how rapidly each oversampler recovers the balanced baseline and whether full recovery is achievable.

The uniform random removal strategy ensures that no distributional criterion biases the removal sequence: each minority instance is equally likely to be removed at each step.
This provides the cleanest possible test of the imbalance hypothesis, isolating pure imbalance effects from any confound introduced by a structured removal policy.


% ============================================================================
\section{Uniform Random Removal Protocol}
\label{sec:causal:random_removal}

The key design choice of Phase~1 is the use of uniform random removal.
At each degradation step, a minority instance $\vx^*$ is selected uniformly at random from the remaining set $\mathcal{X}_{\min}^{(t)}$.

Random removal ensures that the degradation protocol introduces no distributional bias: each instance is equally likely to be removed, providing the cleanest possible test of the imbalance hypothesis.
Unlike structured removal strategies that preferentially target specific regions of the feature space, uniform random removal avoids confounding the effect of reduced count with the effect of a particular removal policy.
This ensures that at each degradation step, the performance drop is attributable to the reduced count (and thus increased $\IR$) and the accompanying information loss, rather than to a systematic bias in which instances were removed.


% ============================================================================
\section{Research Questions}
\label{sec:causal:rqs}

\begin{enumerate}[label=\textbf{RQ\arabic*.}, leftmargin=*, itemsep=4pt]
  \item Does degradation monotonically decrease classification metrics, or is there a threshold imbalance ratio below which performance collapses~\cite{japkowicz2002class, weiss2003effect}?
  \item Are some classifiers more robust to progressive imbalance than others~\cite{ling2004exploration, galar2012review}?
  \item Can oversampling methods fully recover balanced-dataset performance~\cite{chawla2002smote, douzas2018improving}?
  \item Do circular oversampling methods (GVM-CO) recover performance faster or to a higher ceiling than linear interpolation methods (SMOTE, B-SMOTE)?
  \item Is the performance loss caused by imbalance \emph{per se}, or primarily by information loss from a smaller minority sample~\cite{weiss2003effect}?
\end{enumerate}

These five research questions directly connect to the broader thesis narrative.
RQ1--RQ2 characterize the causal signature of imbalance, providing the empirical foundation that motivates Part~IV's oversampling work.
RQ3--RQ4 evaluate whether Part~IV's methods (GVM-CO) succeed at the recovery task.
RQ5 provides the final interpretive synthesis: even when imbalance is the confirmed cause of degradation, irreducible information loss limits the achievable recovery.


